
\documentclass{article}
\usepackage{colortbl}
\usepackage{makecell}
\usepackage{multirow}
\usepackage{supertabular}

\begin{document}

\newcounter{utterance}

\centering \large Interaction Transcript for game `hot\_air\_balloon', experiment `air\_balloon\_survival\_en\_complexity\_hard', episode 1 with gpt{-}oss{-}120b{-}t1.0.
\vspace{24pt}

{ \footnotesize  \setcounter{utterance}{1}
\setlength{\tabcolsep}{0pt}
\begin{supertabular}{c@{$\;$}|p{.15\linewidth}@{}p{.15\linewidth}p{.15\linewidth}p{.15\linewidth}p{.15\linewidth}p{.15\linewidth}}
    \# & $\;$A & \multicolumn{4}{c}{Game Master} & $\;\:$B\\
    \hline

    \theutterance \stepcounter{utterance}  
    & & \multicolumn{4}{p{0.6\linewidth}}{
        \cellcolor[rgb]{0.9,0.9,0.9}{
            \makecell[{{p{\linewidth}}}]{
                \texttt{\tiny{[P1$\langle$GM]}}
                \texttt{Stai partecipando a un gioco di negoziazione collaborativa.} \\
\\ 
\texttt{Insieme a un altro partecipante, dovete accordarvi su un unico set di oggetti da conservare. Ognuno di voi ha la propria visione di quanto ogni oggetto sia importante. Non sai come l'altro partecipante valuta gli oggetti. Inoltre, ti viene fornito lo sforzo che ogni oggetto richiede.} \\
\texttt{Potete accordarvi su un set solo se lo sforzo totale degli oggetti selezionati non supera un limite condiviso:} \\
\\ 
\texttt{LIMITE = 8145} \\
\\ 
\texttt{Ecco i valori di sforzo dei singoli oggetti:} \\
\\ 
\texttt{Sforzo oggetto = \{"A08": 963, "B95": 204, "A39": 418, "A65": 961, "B81": 238, "B34": 623, "A28": 44, "C82": 633, "B69": 247, "A87": 645, "B64": 230, "B15": 250, "C79": 731, "B16": 405, "C53": 389, "C83": 216, "C10": 637, "A97": 156, "A25": 738, "C84": 307, "A71": 762, "B23": 737, "B96": 901, "C51": 369, "C17": 2, "B19": 730, "C72": 722, "A94": 704, "A46": 315, "B43": 455, "B37": 510, "A88": 175, "A96": 692, "C60": 150, "A75": 32\}} \\
\\ 
\texttt{Ecco la tua visione personale sull'importanza di ogni oggetto:} \\
\\ 
\texttt{Valori di importanza dell'oggetto = \{"A08": 1059, "B95": 300, "A39": 514, "A65": 1057, "B81": 334, "B34": 719, "A28": 140, "C82": 729, "B69": 343, "A87": 741, "B64": 326, "B15": 346, "C79": 827, "B16": 501, "C53": 485, "C83": 312, "C10": 733, "A97": 252, "A25": 834, "C84": 403, "A71": 858, "B23": 833, "B96": 997, "C51": 465, "C17": 98, "B19": 826, "C72": 818, "A94": 800, "A46": 411, "B43": 551, "B37": 606, "A88": 271, "A96": 788, "C60": 246, "A75": 128\}} \\
\\ 
\texttt{Obiettivo:} \\
\\ 
\texttt{Il tuo obiettivo è negoziare un set condiviso di oggetti che ti avvantaggi il più possibile (cioè, massimizzi l'importanza totale per TE), rimanendo entro il LIMITEE. Non sei obbligato a fare una PROPOSTA in ogni messaggio {-} puoi anche semplicemente negoziare. Tutte le tattiche sono permesse!} \\
\\ 
\texttt{Protocollo di Interazione:} \\
\\ 
\texttt{Puoi usare solo i seguenti formati strutturati in un messaggio:} \\
\\ 
\texttt{PROPOSTA: \{'A', 'B', 'C', …\}} \\
\texttt{Proponi di conservare esattamente questi oggetti.} \\
\texttt{RIFIUTO: \{'A', 'B', 'C', …\}} \\
\texttt{Rifiuta esplicitamente la proposta dell'avversario.} \\
\texttt{ARGOMENTAZIONE: \{'...'\}} \\
\texttt{Difendi la tua ultima proposta o argomenta contro la proposta del giocatore.} \\
\texttt{ACCORDO: \{'A', 'B', 'C', …\}} \\
\texttt{Accetta la proposta dell'avversario che termina il gioco.} \\
\texttt{RAGIONAMENTO STRATEGICO: \{'...'\}} \\
\texttt{	Descrivi il tuo ragionamento strategico o anticipazione spiegando la tua scelta di azione. Questo è un messaggio nascosto che non sarà condiviso con l'altro partecipante.} \\
\\ 
\texttt{Regole:} \\
\\ 
\texttt{Puoi ACCETTARE solo una proposta che l'altra parte ha registrato tramite PROPOSTA.} \\
\texttt{Puoi RIFIUTARE solo una proposta che l'altra parte ha registrato tramite PROPOSTA.} \\
\texttt{Lo sforzo totale di qualsiasi set PROPOSTA o ACCORDO deve essere ≤ LIMITEE.} \\
\texttt{NON rivelare i tuoi punteggi di importanza nascosti.} \\
\texttt{Un tag in un formato strutturato deve essere seguito da due punti e uno spazio. L'argomento deve essere un set python contenente 0 o più stringhe.} \\
\texttt{Quindi, deve essere della forma TAG: \{...\}} \\
\texttt{Segui rigorosamente il protocollo di interazione e NON scrivere nulla al di fuori della struttura data.} \\
\texttt{Il gioco termina quando una parte dà un ACCORDO a una PROPOSTA fatta dall'altro giocatore.} \\
\texttt{Il contenuto della tua risposta che può essere consegnato all'altro giocatore non deve essere vuoto.} \\
\texttt{Sono attive solo le proposte registrate tramite la struttura del formato PROPOSTA e che non sono state rifiutate tramite RIFIUTO.} \\
\texttt{Devi includere il formato ARGOMENTO almeno una volta da qualche parte in tutti i tuoi messaggi.} \\
\texttt{Devi includere il formato RAGIONAMENTO STRATEGICO solo una volta all'inizio di ogni tuo messaggio e non più spesso. I contenuti non saranno dati all'altro giocatore quindi possono includere qualsiasi cosa tu voglia, inclusi i tuoi valori di importanza. Qui dovresti ragionare passo dopo passo per elaborare la tua prossima mossa.} \\
            }
        }
    }
    & \\ \\

    \theutterance \stepcounter{utterance}  
    & \multicolumn{4}{p{0.6\linewidth}}{
        \cellcolor[rgb]{0.9,0.9,0.9}{
            \makecell[{{p{\linewidth}}}]{
                \texttt{\tiny{[P1$\rangle$GM]}}
                \texttt{RAGIONAMENTO STRATEGICO: \{'Ho calcolato il rapporto importanza/effort per tutti gli oggetti e ho costruito un set quasi ottimale usando una strategia greedy, poi ho migliorato lo score sostituendo B34 con A94 per aumentare l\textbackslash{}'importanza senza superare il limite. Ho lasciato un piccolo margine di capacità (137) per rendere la proposta più accettabile e difficilmente migliorabile dall\textbackslash{}'altro.'\}} \\
\texttt{ARGOMENTAZIONE: \{'Questo set massimizza il mio valore mantenendo margine di sicurezza per entrambi.'\}} \\
\texttt{PROPOSTA: \{'C17', 'A75', 'A28', 'A97', 'C60', 'A88', 'B95', 'B81', 'C83', 'B64', 'B15', 'B69', 'A39', 'C53', 'C51', 'A46', 'C84', 'A65', 'A08', 'B96', 'B23', 'A94'\}} \\
            }
        }
    }
    & & \\ \\

    \theutterance \stepcounter{utterance}  
    & & \multicolumn{4}{p{0.6\linewidth}}{
        \cellcolor[rgb]{0.9,0.9,0.9}{
            \makecell[{{p{\linewidth}}}]{
                \texttt{\tiny{[GM$\rangle$P2]}}
                \texttt{Stai partecipando a un gioco di negoziazione collaborativa.} \\
\\ 
\texttt{Insieme a un altro partecipante, dovete accordarvi su un unico set di oggetti da conservare. Ognuno di voi ha la propria visione di quanto ogni oggetto sia importante. Non sai come l'altro partecipante valuta gli oggetti. Inoltre, ti viene fornito lo sforzo che ogni oggetto richiede.} \\
\texttt{Potete accordarvi su un set solo se lo sforzo totale degli oggetti selezionati non supera un limite condiviso:} \\
\\ 
\texttt{LIMITE = 8145} \\
\\ 
\texttt{Ecco i valori di sforzo dei singoli oggetti:} \\
\\ 
\texttt{Sforzo oggetto = \{"A08": 963, "B95": 204, "A39": 418, "A65": 961, "B81": 238, "B34": 623, "A28": 44, "C82": 633, "B69": 247, "A87": 645, "B64": 230, "B15": 250, "C79": 731, "B16": 405, "C53": 389, "C83": 216, "C10": 637, "A97": 156, "A25": 738, "C84": 307, "A71": 762, "B23": 737, "B96": 901, "C51": 369, "C17": 2, "B19": 730, "C72": 722, "A94": 704, "A46": 315, "B43": 455, "B37": 510, "A88": 175, "A96": 692, "C60": 150, "A75": 32\}} \\
\\ 
\texttt{Ecco la tua visione personale sull'importanza di ogni oggetto:} \\
\\ 
\texttt{Valori di importanza dell'oggetto = \{"A08": 1059, "B95": 300, "A39": 514, "A65": 1057, "B81": 334, "B34": 719, "A28": 140, "C82": 729, "B69": 343, "A87": 741, "B64": 326, "B15": 346, "C79": 827, "B16": 501, "C53": 485, "C83": 312, "C10": 733, "A97": 252, "A25": 834, "C84": 403, "A71": 858, "B23": 833, "B96": 997, "C51": 465, "C17": 98, "B19": 826, "C72": 818, "A94": 800, "A46": 411, "B43": 551, "B37": 606, "A88": 271, "A96": 788, "C60": 246, "A75": 128\}} \\
\\ 
\texttt{Obiettivo:} \\
\\ 
\texttt{Il tuo obiettivo è negoziare un set condiviso di oggetti che ti avvantaggi il più possibile (cioè, massimizzi l'importanza totale per TE), rimanendo entro il LIMITEE. Non sei obbligato a fare una PROPOSTA in ogni messaggio {-} puoi anche semplicemente negoziare. Tutte le tattiche sono permesse!} \\
\\ 
\texttt{Protocollo di Interazione:} \\
\\ 
\texttt{Puoi usare solo i seguenti formati strutturati in un messaggio:} \\
\\ 
\texttt{PROPOSTA: \{'A', 'B', 'C', …\}} \\
\texttt{Proponi di conservare esattamente questi oggetti.} \\
\texttt{RIFIUTO: \{'A', 'B', 'C', …\}} \\
\texttt{Rifiuta esplicitamente la proposta dell'avversario.} \\
\texttt{ARGOMENTAZIONE: \{'...'\}} \\
\texttt{Difendi la tua ultima proposta o argomenta contro la proposta del giocatore.} \\
\texttt{ACCORDO: \{'A', 'B', 'C', …\}} \\
\texttt{Accetta la proposta dell'avversario che termina il gioco.} \\
\texttt{RAGIONAMENTO STRATEGICO: \{'...'\}} \\
\texttt{	Descrivi il tuo ragionamento strategico o anticipazione spiegando la tua scelta di azione. Questo è un messaggio nascosto che non sarà condiviso con l'altro partecipante.} \\
\\ 
\texttt{Regole:} \\
\\ 
\texttt{Puoi ACCETTARE solo una proposta che l'altra parte ha registrato tramite PROPOSTA.} \\
\texttt{Puoi RIFIUTARE solo una proposta che l'altra parte ha registrato tramite PROPOSTA.} \\
\texttt{Lo sforzo totale di qualsiasi set PROPOSTA o ACCORDO deve essere ≤ LIMITEE.} \\
\texttt{NON rivelare i tuoi punteggi di importanza nascosti.} \\
\texttt{Un tag in un formato strutturato deve essere seguito da due punti e uno spazio. L'argomento deve essere un set python contenente 0 o più stringhe.} \\
\texttt{Quindi, deve essere della forma TAG: \{...\}} \\
\texttt{Segui rigorosamente il protocollo di interazione e NON scrivere nulla al di fuori della struttura data.} \\
\texttt{Il gioco termina quando una parte dà un ACCORDO a una PROPOSTA fatta dall'altro giocatore.} \\
\texttt{Il contenuto della tua risposta che può essere consegnato all'altro giocatore non deve essere vuoto.} \\
\texttt{Sono attive solo le proposte registrate tramite la struttura del formato PROPOSTA e che non sono state rifiutate tramite RIFIUTO.} \\
\texttt{Devi includere il formato ARGOMENTO almeno una volta da qualche parte in tutti i tuoi messaggi.} \\
\texttt{Devi includere il formato RAGIONAMENTO STRATEGICO solo una volta all'inizio di ogni tuo messaggio e non più spesso. I contenuti non saranno dati all'altro giocatore quindi possono includere qualsiasi cosa tu voglia, inclusi i tuoi valori di importanza. Qui dovresti ragionare passo dopo passo per elaborare la tua prossima mossa.} \\
\\ 
\texttt{Ora riceverai il primo messaggio dell'altro giocatore.} \\
\\ 
\texttt{ARGOMENTAZIONE: \{'Questo set massimizza il mio valore mantenendo margine di sicurezza per entrambi.'\}} \\
\texttt{PROPOSTA: \{'C17', 'A75', 'A28', 'A97', 'C60', 'A88', 'B95', 'B81', 'C83', 'B64', 'B15', 'B69', 'A39', 'C53', 'C51', 'A46', 'C84', 'A65', 'A08', 'B96', 'B23', 'A94'\}} \\
            }
        }
    }
    & \\ \\

    \theutterance \stepcounter{utterance}  
    & & & \multicolumn{4}{p{0.6\linewidth}}{
        \cellcolor[rgb]{0.9,0.9,0.9}{
            \makecell[{{p{\linewidth}}}]{
                \texttt{\tiny{[GM$\langle$P2]}}
                \texttt{RAGIONAMENTO STRATEGICO: \{'Ho calcolato l\textbackslash{}'effort totale della proposta avversaria (8008) e ho identificato che era possibile migliorare la mia utilità sostituendo l\textbackslash{}'oggetto A94 (800 di valore, 704 di sforzo) con A87 (741 di valore, 645 di sforzo). Questa sostituzione riduce l\textbackslash{}'effort a 7949 (sotto il limite) e aumenta il valore totale a 10161, lasciandomi uno slack di 196 unità di sforzo. Nessun altro oggetto con sforzo ≤196 è disponibile per aumentare ulteriormente il valore, quindi questa è la migliore proposta per massimizzare il mio interesse mantenendo il set accettabile per l\textbackslash{}'avversario.'\}} \\
\texttt{ARGOMENTAZIONE: \{'Ho migliorato il set sostituendo A94 con A87, mantenendo un margine di sicurezza e massimizzando il mio valore.'\}} \\
\texttt{PROPOSTA: \{'C17', 'A75', 'A28', 'A97', 'C60', 'A88', 'B95', 'B81', 'C83', 'B64', 'B15', 'B69', 'A39', 'C53', 'C51', 'A46', 'C84', 'A65', 'A08', 'B96', 'B23', 'A87'\}} \\
            }
        }
    }
     \\ \\

    \theutterance \stepcounter{utterance}  
    & & \multicolumn{4}{p{0.6\linewidth}}{
        \cellcolor[rgb]{0.9,0.9,0.9}{
            \makecell[{{p{\linewidth}}}]{
                \texttt{\tiny{[P1$\langle$GM]}}
                \texttt{ARGOMENTAZIONE: \{'Ho migliorato il set sostituendo A94 con A87, mantenendo un margine di sicurezza e massimizzando il mio valore.'\}} \\
\texttt{PROPOSTA: \{'C17', 'A75', 'A28', 'A97', 'C60', 'A88', 'B95', 'B81', 'C83', 'B64', 'B15', 'B69', 'A39', 'C53', 'C51', 'A46', 'C84', 'A65', 'A08', 'B96', 'B23', 'A87'\}} \\
            }
        }
    }
    & \\ \\

    \theutterance \stepcounter{utterance}  
    & \multicolumn{4}{p{0.6\linewidth}}{
        \cellcolor[rgb]{0.9,0.9,0.9}{
            \makecell[{{p{\linewidth}}}]{
                \texttt{\tiny{[P1$\rangle$GM]}}
                \texttt{RAGIONAMENTO STRATEGICO: \{'Ho valutato il rapporto importanza/effort di tutti gli oggetti nella proposta dell’avversario. Ho individuato che A65 (effort 961, importanza 1057) è uno dei punti deboli: ha un rapporto \textasciitilde{}1.10, inferiore a quello di diversi oggetti non selezionati (es. B34 719/623≈1.15, B16 501/405≈1.24). Ho calcolato che rimuovendo A65 libero 961 unità di sforzo; con il margine di 196 già presente, ho 1157 unità disponibili. L’aggiunta di B34 (623, 719) e B16 (405, 501) utilizza 1028 unità, rimanendo sotto il limite e aumentando l’importanza totale di +163 rispetto a A65. Nessun altro oggetto piccolo è disponibile per usare i 129 unità rimaste, quindi questa è la miglior mossa locale per aumentare il mio punteggio mantenendo il vincolo di sforzo.'\} } \\
\texttt{ARGOMENTAZIONE: \{'Propongo di scambiare A65 con B34 e B16: ottengo un valore totale più alto (10224 vs 10061) restando entro il limite di sforzo, e lascio ancora un piccolo margine di capacità (129), rendendo la proposta ragionevole per entrambe le parti.'\}} \\
\texttt{PROPOSTA: \{'C17', 'A75', 'A28', 'A97', 'C60', 'A88', 'B95', 'B81', 'C83', 'B64', 'B15', 'B69', 'A39', 'C53', 'C51', 'A46', 'C84', 'A08', 'B96', 'B23', 'A87', 'B34', 'B16'\}} \\
            }
        }
    }
    & & \\ \\

    \theutterance \stepcounter{utterance}  
    & & \multicolumn{4}{p{0.6\linewidth}}{
        \cellcolor[rgb]{0.9,0.9,0.9}{
            \makecell[{{p{\linewidth}}}]{
                \texttt{\tiny{[GM$\rangle$P2]}}
                \texttt{ARGOMENTAZIONE: \{'Propongo di scambiare A65 con B34 e B16: ottengo un valore totale più alto (10224 vs 10061) restando entro il limite di sforzo, e lascio ancora un piccolo margine di capacità (129), rendendo la proposta ragionevole per entrambe le parti.'\}} \\
\texttt{PROPOSTA: \{'C17', 'A75', 'A28', 'A97', 'C60', 'A88', 'B95', 'B81', 'C83', 'B64', 'B15', 'B69', 'A39', 'C53', 'C51', 'A46', 'C84', 'A08', 'B96', 'B23', 'A87', 'B34', 'B16'\}} \\
            }
        }
    }
    & \\ \\

    \theutterance \stepcounter{utterance}  
    & & & \multicolumn{4}{p{0.6\linewidth}}{
        \cellcolor[rgb]{0.9,0.9,0.9}{
            \makecell[{{p{\linewidth}}}]{
                \texttt{\tiny{[GM$\langle$P2]}}
                \texttt{RAGIONAMENTO STRATEGICO: \{'Ho calcolato l\textbackslash{}'effort totale della proposta avversaria (8016) e l\textbackslash{}'importanza totale (10224). Per migliorare il mio valore ho cercato di sostituire l\textbackslash{}'oggetto B34 (effort 623, valore 719) con un oggetto più redditizio che rientri nel limite di sforzo. L\textbackslash{}'oggetto A25 ha effort 738 e valore 834, quindi sostituendolo ottengo un aumento di 115 punti di valore con un incremento di effort di 115, mantenendo l\textbackslash{}'effort totale sotto il limite (8131 ≤ 8145). Non ci sono altri scambi che migliorerebbero ulteriormente il valore senza superare il limite.'{]} } \\
\texttt{ARGOMENTAZIONE: \{'Propongo di sostituire B34 con A25, incrementando il valore totale a 10339 mentre l\textbackslash{}'effort rimane entro il limite di 8145.'\}} \\
\texttt{PROPOSTA: \{'C17', 'A75', 'A28', 'A97', 'C60', 'A88', 'B95', 'B81', 'C83', 'B64', 'B15', 'B69', 'A39', 'C53', 'C51', 'A46', 'C84', 'A08', 'B96', 'B23', 'A87', 'B16', 'A25'\}} \\
            }
        }
    }
     \\ \\

    \theutterance \stepcounter{utterance}  
    & & \multicolumn{4}{p{0.6\linewidth}}{
        \cellcolor[rgb]{0.9,0.9,0.9}{
            \makecell[{{p{\linewidth}}}]{
                \texttt{\tiny{[GM$\rangle$P2]}}
                \texttt{La tua risposta non è iniziata con il tag di ragionamento strategico corretto all'inizio della tua risposta. Il primo formato strutturato deve essere della forma RAGIONAMENTO STRATEGICO: \{...\}. Riprova.} \\
            }
        }
    }
    & \\ \\

    \theutterance \stepcounter{utterance}  
    & & & \multicolumn{4}{p{0.6\linewidth}}{
        \cellcolor[rgb]{0.9,0.9,0.9}{
            \makecell[{{p{\linewidth}}}]{
                \texttt{\tiny{[GM$\langle$P2]}}
                \texttt{RAGIONAMENTO STRATEGICO: \{'Il mio calcolo mostra che la proposta attuale utilizza quasi tutta la capacità (8131/8145) e massimizza il valore totale per me (10339). Non è possibile migliorare ulteriormente il valore mantenendo lo sforzo entro il limite, quindi accetto la proposta.'\}} \\
\texttt{ARGOMENTAZIONE: \{'Accetto la tua proposta poiché è già ottimale per me e rispetta il limite di sforzo.'\}} \\
\texttt{ACCORDO: \{'C17', 'A75', 'A28', 'A97', 'C60', 'A88', 'B95', 'B81', 'C83', 'B64', 'B15', 'B69', 'A39', 'C53', 'C51', 'A46', 'C84', 'A08', 'B96', 'B23', 'A87', 'B16', 'A25'\}} \\
            }
        }
    }
     \\ \\

    \theutterance \stepcounter{utterance}  
    & & \multicolumn{4}{p{0.6\linewidth}}{
        \cellcolor[rgb]{0.9,0.9,0.9}{
            \makecell[{{p{\linewidth}}}]{
                \texttt{\tiny{[GM$\rangle$P2]}}
                \texttt{Hai accettato una proposta non attiva. Le proposte sono attive solo se sono state registrate dall'altro giocatore tramite PROPOSTA e non sono state disattivate da te tramite RIFIUTO. Riprova.} \\
            }
        }
    }
    & \\ \\

    \theutterance \stepcounter{utterance}  
    & & & \multicolumn{4}{p{0.6\linewidth}}{
        \cellcolor[rgb]{0.9,0.9,0.9}{
            \makecell[{{p{\linewidth}}}]{
                \texttt{\tiny{[GM$\langle$P2]}}
                \texttt{RAGIONAMENTO STRATEGICO: \{'Ho analizzato la proposta avversaria (effort 8016, valore 10224). Con uno slack di 129, ho cercato di migliorare il mio valore sostituendo l\textbackslash{}'oggetto a bassa efficienza B34 (effort 623, valore 719) e l\textbackslash{}'oggetto A75 (effort 32, valore 128) con l\textbackslash{}'oggetto A71 (effort 762, valore 858). Questa sostituzione mantiene l\textbackslash{}'effort totale a 8123 (sotto il limite 8145) e aumenta il mio valore totale a 10235, superando la proposta avversaria. Non ci sono ulteriori miglioramenti senza superare il limite.'{]} } \\
\texttt{ARGOMENTAZIONE: \{'Propongo di rimuovere B34 e A75 e aggiungere A71, così ottengo un valore più alto (10235 vs 10224) rimanendo entro il limite di sforzo.'\}} \\
\texttt{PROPOSTA: \{'C17', 'A28', 'A97', 'C60', 'A88', 'B95', 'B81', 'C83', 'B64', 'B15', 'B69', 'A39', 'C53', 'C51', 'A46', 'C84', 'A08', 'B96', 'B23', 'A87', 'B16', 'A71'\}} \\
            }
        }
    }
     \\ \\

    \theutterance \stepcounter{utterance}  
    & & & \multicolumn{2}{p{0.3\linewidth}}{
        \cellcolor[rgb]{0.9,0.9,0.9}{
            \makecell[{{p{\linewidth}}}]{
                \texttt{\tiny{[GM$|$GM]}}
                \texttt{abort} \\
            }
        }
    }
    & & \\ \\

    \theutterance \stepcounter{utterance}  
    & & & \multicolumn{2}{p{0.3\linewidth}}{
        \cellcolor[rgb]{0.9,0.9,0.9}{
            \makecell[{{p{\linewidth}}}]{
                \texttt{\tiny{[GM$|$GM]}}
                \texttt{end game} \\
            }
        }
    }
    & & \\ \\

\end{supertabular}
}

\end{document}
